\begin{abstract}
Segmentation quality scores serve different decisions, such as ranking models
and flagging failed predictions, yet are usually reported as if one number
served them all. We test this on eight six-model panels in three domains:
general scenes (COCO, PASCAL VOC, and three ADE20K panels including recent
architectures), street scenes (Cityscapes, KITTI), and abdominal CT (AMOS). For
seven declared criteria, every scalar ordering of the models must contradict
\NBAllMin{}--\NBAllMax{} of the non-tied criterion preferences, and exhaustive
search over all 720 orderings shows this bound is essentially exact. Scoring
each declared endpoint with its own frozen formula instead of one mixed score
improves policy-rank agreement, operational-rank agreement, and severe-miss
AUROC on all eight panels; all \NNontrivialSig{} contrasts are significant after
Holm correction over \NFamilyN{} tests. A single failure score is more fragile.
Selected on general scenes, it beats the strongest single metric there, but not
on street scenes and CT, where high-error failures make up at most
\NMaxShareLose\% of failures and severe misses dominate. Re-selecting it on
Cityscapes training data brings it level with the strongest single metric on
three held-out panels, including CT (gains of \NSelGainCity, \NSelGainKitti, and
\NSelGainAmos{} AUROC over the general-scene selection). Quality
evaluation should therefore name its endpoint and match its failure score to
the failure composition of the target domain.
\end{abstract}
